\section{Quantum}

Here is the hardest test the base faces. It is local, discrete, and deterministic, each dock reading only its own sites, the beat a fixed table, nothing reaching across the mesh. Quantum mechanics is none of these, and two distant measurements on a shared past show correlations no local story can match. Rather than bury the tension, the paper measures it.

The measure is the CHSH correlation. Any local deterministic theory, any theory where outcomes are fixed in advance by what each side carries, is capped at $2$. Quantum mechanics reaches $2\sqrt 2$, the Tsirelson bound, and experiment agrees with quantum mechanics. Worked on the substrate, with the measurement itself modeled as part of the mesh rather than imposed from outside, the correlation climbs to the quantum value, and it does so only when the two sides share an aligned past, decaying back to the classical bound as that past is scrambled.

\begin{result}[The quantum correlations]
On the substrate the CHSH correlation reaches the Tsirelson bound $2\sqrt 2$ for an aligned shared past and falls to the classical $2$ for a generic one. The Born rule's square law is forced by additivity, and a quantum walk spreads ballistically where a classical one only diffuses. \cite{pollard2026vibetest}
\end{result}

More of the quantum formalism follows. The Born rule, the square that turns an amplitude into a probability, is forced once you ask for a law additive over alternatives, and the exponent comes out two. A quantum walk on the mesh spreads at constant speed where a classical random walk only crawls as the square root of time, the ballistic signature of genuine quantum motion. A pair drawn together by an attractive coupling settles into a localized bound state well below the free band, with discrete internal levels, a hydrogen-like spectrum on the lattice. And the least classical formulation of all, a sum over histories run on a discrete causal set, recovers a spacetime of about the right dimension, near three in the full construction, so even the path integral lands on the geometry the rest of the paper built. What stays open is the measurement itself, the collapse, derived rather than assumed. The Bell number is reached and the Born rule is forced, but a full account of measurement is still partial, and that is the frontier where the theory is most exposed and most interesting at once.

The reading the model favors is that a measurement is not a collapse but a settling. A coherent body does not flip the instant new influences reach it. Its inner tones must re-reach agreement before it emits a fresh collective tone, and that interval, felt from the inside as weighing the options, is from the outside a plain deterministic relaxation into one definite state. The single classical record appears because the fine phase leaks into the bath and the body's own boundary screens its inside from its outside, so the coarse description sees one outcome, with no special collapse law added. What a full reconstruction still owes is a derivation of this from first principles, and the standard obligations any such route must meet, that the underlying values be read relationally rather than fixed, and that the crystal's loops force the phase to wind in whole turns.

There is a marked route toward closing the gap, borrowed rather than invented. Seen the right way the knit is simpler than its table suggests. The ternary tone is a single spin, with its three values the three settings of that spin, and the collision's three telling moves, create, annihilate, and hop, are one and the same operation, the transfer of one unit of tone across a line. So the knit is a charge-conserving exchange of tone quanta, and for exactly that class a standard construction rewrites the dynamics, with no approximation, as a field theory with creation and annihilation operators, handing over the field and its action in one step.

What that step yields is at first a lifeless, Euclidean theory, and whether it is genuinely quantum, with a true Hilbert space and the right probabilities, is decided by a single checkable property of the lattice weights, reflection positivity, the theorem that licenses the turn from a statistical model to a unitary quantum one. Here the reversibility already proven does real work, because it is precisely what turns the otherwise lifeless generator into a Hermitian, genuinely quantum Hamiltonian, and the imaginary unit appears with it, as the geometric phase the tone sweeps as it turns. The continuum limit of that Hamiltonian is a free field carrying the conserved charge, a thing already known to be reflection positive, so the test is expected to pass once taken near its critical point, and confirming it is the decisive next experiment. None of this is settled. It is the path from the substrate to a rigorous quantum field, and it is built to be walked by computation.
